\documentclass[11pt,a4paper]{article}

% --- Pacotes de Idioma e Codificação ---
\usepackage[utf8]{inputenc}
\usepackage[T1]{fontenc}
\usepackage[portuguese]{babel}

% --- Design e Layout ---
\usepackage[margin=2.5cm]{geometry}
\usepackage{charter} % Fonte profissional e legível
\usepackage{lastpage}
\usepackage{fancyhdr}
\usepackage[table]{xcolor}
\usepackage{tabularx}
\usepackage{booktabs}
\usepackage{xcolor}
\usepackage{enumitem}
\usepackage{amssymb}

% --- Configuração de Cores ---
\definecolor{primaryblue}{RGB}{0, 51, 102}

% --- Cabeçalho e Rodapé ---
\pagestyle{fancy}
\fancyhf{}
\fancyhead[L]{\small \textbf{Registo de Entrevista} | Projeto Interno}
\fancyhead[R]{\small \thepage\ de \pageref{LastPage}}
\fancyfoot[L]{\small \textit{Confidencial - Apenas para uso interno}}
\renewcommand{\headrulewidth}{0.4pt}

% --- Ambiente Personalizado para Transcrição ---
% Facilita a distinção entre Entrevistador (E) e Entrevistado (R)
\newlist{transcription}{description}{1}
\setlist[transcription]{
	labelsep=1em,
	leftmargin=3.5em,
	style=nextline,
	font=\bfseries\color{primaryblue}
}

\newcommand{\pergunta}[1]{\item[E:] \textit{#1}}
\newcommand{\resposta}[1]{\item[R:] #1}

% --- Início do Documento ---
\begin{document}
	
	% Título Principal
	{\noindent\LARGE\bfseries\color{primaryblue} Relatório de Entrevista Individual}
	\\[0.5cm]
	
	% --- Tabela de Metadados ---
	\renewcommand{\arraystretch}{1.5}
	\begin{tabularx}{\textwidth}{@{}|l|X|l|X|@{}}
		\hline
		\rowcolor{gray!10} \textbf{ID Entrevista} & \#003 & \textbf{Data} & 26 de Fevereiro de 2026 \\ \hline
		\textbf{Hora Início} & 14:00 & \textbf{Hora Fim} & 14:40 \\ \hline
		\textbf{Local/Meio} & Sala de Reuniões B & \textbf{Estado} & Finalizada / Transcrita \\ \hline
	\end{tabularx}
	
	\vspace{0.5cm}
	
	% --- Participantes ---
	\begin{tabularx}{\textwidth}{@{}|l|X|@{}}
		\hline
		\rowcolor{gray!10} \textbf{Intervenientes} & \textbf{Cargo / Função} \\ \hline
		\textbf{Entrevistador(a)} & [O Teu Nome] - Consultor de Processos \\ \hline
		\textbf{Entrevistado(a)}  & [Nome do Colaborador] - Gestor de Stock \\ \hline
	\end{tabularx}
	
	\vspace{0.8cm}
	
	% --- Contexto e Objetivos ---
	\section*{1. Enquadramento e Objetivos}
	\begin{itemize}[leftmargin=*, label=\small\color{primaryblue}$\blacksquare$]
		\item \textbf{Fidelidade do Inventário}: O objetivo é garantir que o stock "teórico" no computador seja sempre igual ao stock "real" na prateleira, evitando paragens na oficina por falta de material.
		\item \textbf{Rastreabilidade Total}: Garantir que cada parafuso ou motor tem um destino (uma venda ou uma reparação) e um custo associado, protegendo a margem de lucro.
		\item \textbf{Eficiência na Reposição}: Automatizar a identificação de faltas para que o processo de compra seja baseado em dados reais de consumo e não em "palpites".
	\end{itemize}
	
	\hrule % Linha separadora opcional
	\vspace{0.5cm}
	
	% --- Transcrição ---
	\section*{2. Transcrição / Notas da Entrevista}
	
	\begin{transcription}
		
		\pergunta{Quando recebe uma encomenda de um fornecedor, como é que dá entrada das peças no sistema hoje? Confere uma por uma ou introduz o total da fatura?}
		\resposta{Quando recebo uma encomenda de peças de um fornecedor, costumo conferir uma a uma, comparando com a guia de remessa ou a fatura. Verifico se a quantidade e o tipo de peças estão corretos e se não há nada danificado. Depois disso, registo a entrada no nosso caderno ou numa folha de cálculo, onde tenho uma lista com as peças, as quantidades e, por vezes, o preço de compra. Não costumo introduzir só o total da fatura, porque preciso de saber exatamente o que entrou, até para controlar o stock e perceber o que está a sair mais.}
		
		\pergunta{Como identifica as peças? Usa as referências do fornecedor ou prefere criar códigos internos da oficina (ex: 'PNEU-85-XIA')?}
		\resposta{Para identificar as peças, uso uma mistura. Quando o fornecedor já tem uma referência clara e fácil de entender, aproveito essa. Mas, na maioria das vezes, prefiro criar um código interno mais simples, algo como “PNEU-85-XIA” ou “TRAVAO-M365”, porque assim é mais fácil para mim e para os mecânicos encontrarem a peça certa. Esses códigos também ajudam a manter a organização no armário ou nas prateleiras.}
		
		\pergunta{Trabalha com números de série para peças caras (ex: baterias ou motores) para saber exatamente qual a unidade que foi montada em que trotinete?}
		\resposta{No caso de peças mais caras, como baterias ou motores, sim, tento sempre apontar o número de série. Isso é importante para sabermos exatamente qual foi a unidade instalada em cada trotinete, principalmente se houver uma reclamação ou se for preciso acionar a garantia. Nem sempre é fácil, porque às vezes os números vêm mal visíveis ou em etiquetas que se soltam, mas faço o possível para manter esse controlo.}
		
		\pergunta{Como é que uma peça 'sai' do stock para uma reparação? O mecânico pede-lhe a peça e você regista-a logo na Ordem de Serviço, ou o mecânico tira-a e você regista-a mais tarde?}
		\resposta{Quando uma peça é usada numa reparação, o normal é o mecânico vir ter comigo e pedir a peça. Eu entrego e registo logo na folha da ordem de serviço, para garantir que não me esqueço. Mas há dias mais complicados em que o mecânico vai buscar diretamente e só me avisa depois. Nesses casos, tento confirmar o mais depressa possível e atualizar o registo, mas admito que às vezes pode escapar alguma coisa.}
		
		\pergunta{E se um mecânico levantar uma peça para um diagnóstico e depois não a usar? Como é que faz o processo de 'devolução' da peça ao stock?}
		\resposta{Se o mecânico levantar uma peça para testar durante o diagnóstico e depois não a usar, costumo pedir que me devolva logo. Quando isso acontece, volto a colocá-la no sítio e atualizo o stock, como se fosse uma entrada. Tento manter tudo equilibrado, mas seria bom ter uma forma mais prática de fazer esse “vai e volta” das peças, porque hoje em dia é tudo manual e depende muito da nossa memória.}
		
		\pergunta{Como é que sabe que está na hora de encomendar mais material? Gostaria que o sistema tivesse um 'Stock Mínimo' que acendesse um alerta vermelho quando restassem apenas 3 ou 4 unidades?}
		\resposta{Quando preciso de saber se está na hora de encomendar mais material, costumo verificar o stock manualmente ou através das anotações que mantenho atualizadas. Tenho uma ideia dos mínimos que gosto de ter para cada tipo de peça, mas nem sempre é fácil controlar tudo ao detalhe. Por isso, sim, seria muito útil se o sistema tivesse um campo de “stock mínimo” e me avisasse automaticamente quando uma peça estivesse a chegar ao limite, por exemplo, com um alerta vermelho quando só restassem 3 ou 4 unidades. Isso ajudava-me a agir antes de ficar sem material.}
		
		\pergunta{Costuma fazer inventários (contagens) totais ou prefere ir contando família a família (ex: hoje conto pneus, amanhã conto luzes)?}
		\resposta{Quanto aos inventários, não costumo fazer contagens totais com muita frequência porque dá muito trabalho e exige parar parte da oficina. Prefiro ir fazendo por grupos de peças — hoje conto os pneus, amanhã as luzes, depois os travões, e assim por diante. Tento fazer isso todas as semanas, para manter o controlo sem parar tudo de uma vez.}
		
		\pergunta{Como regista as quebras? (ex: uma câmara de ar que se furou na montagem ou uma peça que desapareceu).}
		\resposta{Quando há quebras, como uma câmara de ar que se fura durante a montagem ou uma peça que desaparece, costumo apontar num caderno ou numa folha à parte. Escrevo o que aconteceu, a data e, se souber, quem estava a usar a peça. Depois dou baixa manualmente no stock. Mas seria melhor se o sistema tivesse uma forma de registar essas quebras diretamente, com um motivo associado, para ficar tudo mais organizado e transparente.}
		
		\pergunta{A oficina tem várias prateleiras ou zonas? Seria útil o sistema dizer-lhe exatamente em que corredor ou gaveta está guardada a peça que o mecânico pediu?}
		\resposta{A oficina tem várias prateleiras e zonas, sim. Tenho uma parte para pneus, outra para peças elétricas, outra para consumíveis, e assim por diante. Às vezes perco tempo à procura de uma peça, principalmente se for algo que não uso com frequência. Por isso, sim, seria muito útil se o sistema me dissesse exatamente onde está guardada cada peça — por exemplo, “Gaveta 3, Prateleira B”. Isso tornava tudo mais rápido, tanto para mim como para os mecânicos.}
		
		\pergunta{Existem kits de peças? (ex: 'Kit Revisão' que inclui pastilhas + limpeza). Gostaria de dar saída de um kit e o sistema descontar automaticamente todos os componentes que o compõem?}
		\resposta{Temos alguns kits, sim. Por exemplo, o “Kit de Revisão” que inclui pastilhas, lubrificante e um pano de limpeza. Quando sai um kit, tenho de dar baixa manualmente de cada peça que o compõe, o que dá trabalho e pode levar a erros. Era ótimo se o sistema permitisse dar saída de um kit e ele próprio descontasse automaticamente todas as peças incluídas.}
		
		\pergunta{Guarda os contactos dos fornecedores no sistema? Seria útil conseguir gerar uma lista de compras automática para enviar por e-mail ao fornecedor?}
		\resposta{Guardo os contactos dos fornecedores num ficheiro Excel, com os nomes, telefones e e-mails. Também tenho lá os preços habituais e os prazos de entrega. Se o sistema permitisse gerar uma lista de compras com base no que está em falta e depois enviar isso diretamente por e-mail ao fornecedor, poupava-me muito tempo e evitava esquecimentos.}
		
		\pergunta{Como controla as peças que enviou para trás por estarem com defeito (garantias)? Precisa de um estado no sistema para 'Peças em Reclamação'?}
		\resposta{Quando tenho de devolver peças com defeito, costumo separá-las numa caixa marcada como “para devolução” e aponto num papel o que está lá dentro, com a data e o motivo. Mas às vezes perco o controlo, principalmente se o fornecedor demorar a responder. Seria muito útil se o sistema tivesse um estado específico para “Peças em Reclamação”, onde eu pudesse ver o que está pendente, há quanto tempo, e com que fornecedor.}
		
		
	\end{transcription}
	
%	\vspace{0.8cm}
%	
%	% --- Observações e Próximos Passos ---
%	\section*{3. Análise e Observações Gerais}
%	\begin{itemize}[leftmargin=*]
%		\item \textbf{Tom de voz:} O entrevistado demonstrou abertura, mas alguma frustração com os sistemas de IT atuais.
%		\item \textbf{Pontos Críticos:} Falta de interoperabilidade entre sistemas.
%	\end{itemize}
%	
%	\section*{4. Próximos Passos}
%	\begin{tabularx}{\textwidth}{@{}|X|l|l|@{}}
%		\hline
%		\rowcolor{gray!10} \textbf{Ação} & \textbf{Responsável} & \textbf{Prazo} \\ \hline
%		Validar fluxo de dados com equipa de IT & [Teu Nome] & 05/03/2026 \\ \hline
%		Agendar follow-up para demonstração de solução & [Teu Nome] & 12/03/2026 \\ \hline
%	\end{tabularx}

\end{document}